\centerline{\bf \Huge Конец Евангелиона}


\begin{flushright}
        \scriptsize{Эпизод 25. Ты любишь меня?}
\end{flushright}


Я думал написать здесь какую-то душераздирающую историю о том как сильно на меня повлияла Ева, но я решил что хватит \href{https://www.youtube.com/watch?v=5YY7xEH5FGc}{и этого...}


\problem{1}\textbf{Теорема о существовании базиса.} В любом конечномерном пространстве существует базис.


\problem{2}\textbf{Однозначность размера базиса.} Докажите, что

a) Базис наименьшая по включению полная система векторов

б) Базис это наибольшая по включению линейно независимая система векторов

в)


\problem{3}
 \textbf{Теорема.} Имеют место следующие утверждения:

\begin{enumerate}

    \item[a.]
    система векторов $\left\{X_1, \ldots, X_k\right\}$ с линейно зависимой подсистемой сама линейно зависима;

    \item[b.]
    любая часть линейно независимой системы векторов $\left\{X_1, \ldots, X_k\right\}$ линейно независима;

    \item[c.]
    среди линейно зависимой системы векторов хотя бы один является линейной комбинацией остальных;

    \item[d.]
    если один из векторов системы выражается через остальные, то система линейно зависима;

    \item[e.]
    если векторы $X_1, \ldots, X_k$ линейно независимы, а векторы $X_1, \ldots, X_k, X$ линейно зависимы, то $X$ - линейная комбинация векторов $X_1, \ldots, X_k$;

    \item[f.]
    если векторы $X_1, \ldots, X_k$ линейно независимы $u$ вектор $X_{k+1}$ нельзя через них выразить, то система $X_1, \ldots, X_k, X_{k+1}$ линейно независима.

\end{enumerate}

\problem{4} Докажите, что операция композиции функций ассоциативна.


\newpage


\begin{flushright}
        \scriptsize{Эпизод 26. Береги себя.}
\end{flushright}

\problem{5}Является ли множество группой относительно указанной операции? Если да, является ли эта группа абелевой?

\begin{enumerate}

    \item \((\mathbb Q\setminus\{0\},\cdot)\).

    \item Множество чётных целых чисел относительно сложения.

    \item Множество нечётных целых чисел относительно сложения.

    \item \(\mathbb Z_p\setminus\{0\}\) относительно умножения по модулю \(5\).

    \item \(\mathbb Z_n\setminus\{0\}\) относительно умножения по модулю \(6\). Подумайте, что нужно убрать чтобы это стало группой?

    \item \(\{1,3,5,7\}\) относительно умножения по модулю \(8\).

    \item \(\{-1,1\}\) относительно умножения.


    \item \(\mathbb R\setminus\{-1\}\) относительно операции
    \[
    a*b=a+b+ab.
    \]

    \item Множество всех ненулевых векторов плоскости относительно сложения.

    \item Множество всех поворотов правильного треугольника относительно композиции.

    \item Множество всех симметрий правильного треугольника относительно композиции.
\end{enumerate}


\problem{6} \textbf{Куча задач посложнее}
\begin{enumerate}
    \item Докажите, что в любой группе нейтральный элемент единственен.

    \item Докажите, что в любой группе обратный элемент к данному элементу единственен.

    \item Докажите, что если \(a^2=e\) для любого элемента \(a\) группы \(G\), то группа \(G\) абелева.

    \item Пусть \(G\) --- группа, а \(a,b\in G\). Докажите, что
    \[
        (ab)^{-1}=b^{-1}a^{-1}.
    \]

    \item Пусть \(G\) --- конечная группа, а \(a\in G\). Докажите, что существует такое натуральное \(n\), что
    \[
        a^n=e.
    \]

    \item Найдите порядки всех элементов группы \(\mathbb Z_{12}\) относительно сложения по модулю \(12\).

    \item Докажите, что группа \(\mathbb Z_n\) относительно сложения является циклической.

    \item Найдите все образующие элементы группы \(\mathbb Z_{18}\) относительно сложения по модулю \(18\).

    \item Рассмотрим группу перестановок \(S_3\). Выпишите все её элементы.

    \item В группе \(S_3\) найдите порядки всех элементов.

    \item Докажите, что группа \(S_3\) не является абелевой.

    \item Найдите порядок элемента
    \[
        (1\ 2\ 3)(4\ 5)
    \]
    в группе \(S_5\).

    \item Найдите порядок перестановки
    \[
        (1\ 4\ 2)(3\ 5\ 6\ 7)
    \]
    в группе \(S_7\).

    \item Докажите, что порядок перестановки равен наименьшему общему кратному длин её независимых циклов.

    \item Приведите пример неабелевой группы, в которой не все элементы имеют порядок \(2\).

    \item Приведите пример неабелевой группы, в которой есть элемент порядка \(2\).

    \item Докажите, что группа всех движений правильного \(n\)-угольника имеет \(2n\) элементов.

    \item Докажите, что группа симметрий правильного \(n\)-угольника порождается поворотом \(r\) и отражением \(s\), для которых
    \[
        r^n=e,\qquad s^2=e,\qquad srs=r^{-1}.
    \]

    \item Для группы диэдра \(D_4\) найдите порядок каждого элемента.

    \item Докажите, что группа
    \[
        \mathbb Q\setminus\{0\}
    \]
    относительно умножения не является циклической.

    \item Докажите, что группа
    \[
        \mathbb R\setminus\{0\}
    \]
    относительно умножения не является циклической.
    
\end{enumerate} 